\section{Gambar: Elemen img dan figure}

Gambar merupakan elemen visual yang esensial dalam desain web modern, digunakan untuk menyampaikan informasi, meningkatkan estetika, dan memperkaya pengalaman pengguna. Elemen \code{<img>} adalah elemen kosong (void element) yang menyisipkan gambar ke dalam halaman. Atribut \code{src} (source) menentukan lokasi file gambar, yang dapat berupa URL absolut, path relatif, atau bahasa data (data URI). Atribut \code{alt} (alternative text) menyediakan deskripsi tekstual gambar yang ditampilkan saat gambar gagal dimuat dan dibaca oleh pembaca layar untuk pengguna dengan disabilitas visual. Atribut \code{alt} secara teknis opsional dalam HTML, namun dari perspektif aksesibilitas dan SEO, ia dianggap wajib untuk setiap gambar yang memiliki makna konten \cite{mdnAccessibility}.

Atribut dimensi \code{width} dan \code{height} pada elemen \code{<img>} memiliki fungsi penting yang melampaui sekadar mengatur ukuran tampilan. Ketika browser memuat halaman, ia perlu mengetahui berapa ruang yang harus dialokasikan untuk gambar sebelum gambar sebenarnya selesai diunduh. Tanpa atribut dimensi, browser mengalokasikan ruang minimal dan kemudian melakukan reflow (perhitungan ulang layout) saat gambar dimuat, menyebabkan \textit{layout shift} yang mengganggu pengalaman pengguna - konten dapat tiba-tiba bergerak saat pengguna sedang membaca atau berinteraksi. Menyediakan dimensi memungkinkan browser mereservasi ruang yang tepat, menciptakan pengalaman yang lebih mulus.

HTML5 memperkenalkan elemen \code{<figure>} dan \code{<figcaption>} untuk mengelompokkan konten yang mandiri dengan keterangan yang sesuai. Elemen \code{<figure>} dapat berisi gambar, diagram, cuplikan kode, atau konten mandiri lainnya yang dirujuk dari aliran dokumen utama. Elemen \code{<figcaption>} berisi keterangan atau penjelasan untuk konten dalam \code{<figure>}. Kombinasi ini memberikan semantik yang jelas bahwa keterangan terkait dengan konten spesifik di dalamnya, bukan sekadar teks yang kebetulan berada di dekat gambar. Struktur \code{<figure>} dengan \code{<figcaption>} meningkatkan aksesibilitas karena pembaca layar dapat mengasosiasikan keterangan dengan konten secara eksplisit \cite{webdevAccessibility}.

Format gambar web memiliki karakteristik berbeda yang mempengaruhi kualitas, ukuran file, dan kasus penggunaan. \textbf{JPEG} ideal untuk foto dengan banyak warna dan gradasi, menawarkan kompresi lossy yang efisien. \textbf{PNG} mendukung transparansi dan cocok untuk grafik dengan area warna solid, logo, dan screenshot. \textbf{SVG} (Scalable Vector Graphics) adalah format vektor yang dapat diskalakan tanpa kehilangan kualitas, sempurna untuk ikon, logo, dan diagram. \textbf{WebP} adalah format modern yang menawarkan kompresi superior dibanding JPEG dan PNG dengan dukungan transparansi dan animasi, didukung oleh browser modern. Memilih format yang tepat dapat secara dramatis mengurangi ukuran file dan meningkatkan waktu muat halaman.

\begin{table}[htbp]
\centering
\begin{tabular}{|P{2.5cm}|P{2.5cm}|P{3cm}|P{4cm}|}
\hline
\textbf{Atribut} & \textbf{Wajib?} & \textbf{Tipe} & \textbf{Fungsi} \\
\hline
\code{src} & Ya & URL & Lokasi file gambar \\
\hline
\code{alt} & Rekomendasi kuat & Teks & Deskripsi untuk aksesibilitas \\
\hline
\code{width} & Opsional & Angka (px) & Lebar tampilan \\
\hline
\code{height} & Opsional & Angka (px) & Tinggi tampilan \\
\hline
\code{loading} & Opsional & \code{lazy}, \code{eager} & Strategi loading \\
\hline
\end{tabular}
\caption{Atribut Elemen Gambar (\code{<img>})}
\label{tab:img-attributes}
\end{table}

\begin{webcode}[language=HTML, caption={Contoh Penggunaan Gambar}]
<!-- Gambar dasar dengan alt text -->
<img src="images/foto-profil.jpg" 
     alt="Foto profil John Doe tersenyum di kantor" 
     width="200" height="200" />

<!-- Gambar dengan figure dan figcaption -->
<figure>
  <img src="diagram/alur-proses.png" 
       alt="Diagram alur proses pendaftaran: 
            mulai dari form, verifikasi, sampai aktivasi" 
       width="600" height="400" />
  <figcaption>
    Gambar 1: Diagram alur proses pendaftaran pengguna baru
  </figcaption>
</figure>

<!-- Gambar dengan lazy loading -->
<img src="images/galeri/banyak.jpg" 
     alt="Galeri foto acara" 
     loading="lazy" 
     width="800" height="600" />

<!-- Gambar dekoratif (alt kosong) -->
<img src="images/dekorasi.png" 
     alt="" 
     role="presentation" />
\end{webcode}

\begin{catatan}
\textbf{Praktik Aksesibilitas Gambar:}
\begin{itemize}[leftmargin=*]
  \item Selalu gunakan \code{alt} yang deskriptif untuk gambar bermakna.
  \item Gunakan \code{alt=""} untuk gambar murni dekoratif (tambahkan \code{role="presentation"}).
  \item Jangan gunakan kata "gambar" atau "foto" di alt text - screen reader sudah mengumumkan ini.
  \item Deskripsikan informasi yang relevan, bukan setiap detail visual.
  \item Untuk grafik kompleks, pertimbangkan deskripsi panjang dengan \code{longdesc} atau \code{aria-describedby}.
\end{itemize}
\end{catatan}

